
\documentclass[11pt,letterpaper,oneside]{report}
\usepackage[T1]{fontenc}
\usepackage[utf8]{inputenc}
\usepackage{lmodern}
\usepackage{inconsolata}
\usepackage{geometry}
\geometry{margin=0.89in,headheight=16pt,headsep=19pt,footskip=26pt}
\usepackage{xcolor}
\definecolor{navy}{HTML}{173758}
\definecolor{teal}{HTML}{157B89}
\definecolor{pale}{HTML}{E9F0F3}
\definecolor{codegray}{HTML}{F5F7F9}
\definecolor{muted}{HTML}{667785}
\usepackage{booktabs,longtable,tabularx,array,calc}
\usepackage{amsmath,amssymb}
\usepackage{graphicx}
\usepackage{microtype}
\usepackage{enumitem}
\usepackage{fancyhdr}
\usepackage{titlesec}
\usepackage{listings}
\usepackage{upquote}
\usepackage{xurl}
\usepackage{hyperref}
\hypersetup{colorlinks=true,linkcolor=navy,urlcolor=teal,bookmarksopen=true,pdfauthor={Software Build Engineering Documentation Suite}}
\urlstyle{same}
\setcounter{tocdepth}{2}
\setcounter{secnumdepth}{2}
\setlength{\parindent}{0pt}
\setlength{\parskip}{5pt plus 1pt}
\setlength{\emergencystretch}{3em}
\setlist[itemize]{leftmargin=*,topsep=4pt,itemsep=2pt}
\setlist[enumerate]{leftmargin=*,topsep=4pt,itemsep=2pt}
\titleformat{\chapter}[display]{\normalfont\bfseries\color{navy}}{\fontsize{11}{12}\selectfont CHAPTER \thechapter}{7pt}{\fontsize{23}{28}\selectfont}
\titlespacing*{\chapter}{0pt}{-6pt}{19pt}
\titleformat{\section}{\Large\bfseries\color{navy}}{\thesection}{0.65em}{}
\titleformat{\subsection}{\large\bfseries\color{teal}}{\thesubsection}{0.65em}{}
\pagestyle{fancy}
\fancyhf{}
\fancyhead[L]{\small\textcolor{muted}{PROFESSIONAL CMAKE ENGINEERING}}
\fancyhead[R]{\small\textcolor{muted}{\nouppercase{\leftmark}}}
\fancyfoot[C]{\small\textcolor{muted}{\thepage}}
\renewcommand{\headrulewidth}{0.3pt}
\fancypagestyle{plain}{\fancyhf{}\fancyfoot[C]{\small\thepage}\renewcommand{\headrulewidth}{0pt}}
\lstdefinelanguage{cmake}{sensitive=false,morekeywords={cmake_minimum_required,project,add_library,add_executable,target_compile_features,target_include_directories,target_link_libraries,target_sources,target_compile_options,target_compile_definitions,target_link_options,add_subdirectory,include,if,else,elseif,endif,foreach,endforeach,function,endfunction,macro,endmacro,set,option,find_package,find_program,add_custom_command,add_custom_target,install,configure_file,configure_package_config_file,write_basic_package_version_file,set_target_properties,set_tests_properties,add_test,message,FetchContent_Declare,FetchContent_MakeAvailable,return,cmake_parse_arguments},morecomment=[l]{\#},morestring=[b]"}
\lstdefinelanguage{sh}{morekeywords={cmake,ctest,cpack,find,sha256sum,clang-tidy,echo,cp,make,ninja,git},morecomment=[l]{\#},morestring=[b]"}
\lstdefinelanguage{yaml}{morecomment=[l]{\#},morestring=[b]",morekeywords={name,on,permissions,jobs,steps,uses,run,with,env,matrix,include,if}}
\lstdefinelanguage{text}{}
\lstset{basicstyle=\ttfamily\footnotesize,keywordstyle=\color{navy}\bfseries,commentstyle=\color{muted},stringstyle=\color{teal},backgroundcolor=\color{codegray},frame=single,rulecolor=\color{pale},breaklines=true,breakatwhitespace=false,showstringspaces=false,columns=fullflexible,keepspaces=true,tabsize=2,xleftmargin=5pt,xrightmargin=5pt,aboveskip=9pt,belowskip=9pt}
\providecommand{\tightlist}{\setlength{\itemsep}{0pt}\setlength{\parskip}{0pt}}
\providecommand{\passthrough}[1]{#1}
\providecommand{\real}[1]{#1}

\begin{document}
\hypersetup{pageanchor=false}\begin{titlepage}
\color{navy}
\vspace*{1.15in}
{\Large\bfseries PROFESSIONAL CMAKE\\[7pt]SOFTWARE BUILD ENGINEERING\par}
\vspace{0.32in}\noindent\textcolor{teal}{\rule{0.9\linewidth}{2.3pt}}\par
\vspace{0.45in}
{\parbox{0.95\textwidth}{\raggedright\fontsize{24}{29}\selectfont\bfseries Cross-Platform and Performance Engineering\par}}
\vspace{0.3in}
{\Large\color{teal}Toolchains, portability, measurable optimization\par}
\vfill
{\large DOCUMENT 07 OF 08\par}
\vspace{0.1in}
{\small Companion synthesis to Craig Scott, \emph{Professional CMake: A Practical Guide}, 20th Edition (2025).\\
Original learning/engineering material. Not affiliated with or endorsed by the author.\par}
\vspace{0.2in}
{\small 17 October 2026\par}
\end{titlepage}
\hypersetup{pageanchor=true}\clearpage\pagenumbering{roman}\tableofcontents\clearpage\pagenumbering{arabic}

\textbf{Purpose.} Select and validate portable toolchains,
cross-compilation strategies and measurable build optimizations.
\textbf{{[}BOOK-GROUNDED{]}} Chapters 15--17, 23--26, 28, 35 and 43
explain foundational CMake/toolchain concepts. \textbf{{[}ENGINEERING
EXTENSION{]}} detailed benchmark methodology, platform validation SLOs
and risk matrices are newly proposed methods.

\chapter{Cross-platform build matrix}\label{cross-platform-build-matrix}

\begin{longtable}[]{@{}
  >{\raggedright\arraybackslash}p{(\columnwidth - 6\tabcolsep) * \real{0.2500}}
  >{\raggedright\arraybackslash}p{(\columnwidth - 6\tabcolsep) * \real{0.2500}}
  >{\raggedright\arraybackslash}p{(\columnwidth - 6\tabcolsep) * \real{0.2500}}
  >{\raggedright\arraybackslash}p{(\columnwidth - 6\tabcolsep) * \real{0.2500}}@{}}
\toprule\noalign{}
\begin{minipage}[b]{\linewidth}\raggedright
Platform
\end{minipage} & \begin{minipage}[b]{\linewidth}\raggedright
Common generator
\end{minipage} & \begin{minipage}[b]{\linewidth}\raggedright
Compiler family
\end{minipage} & \begin{minipage}[b]{\linewidth}\raggedright
Critical deployment difference
\end{minipage} \\
\midrule\noalign{}
\endhead
\bottomrule\noalign{}
\endlastfoot
Linux x86-64 & Ninja / Unix Makefiles & GCC / Clang & RPATH, shared
object loader and standard filesystem layout \\
Windows x64 & Visual Studio / Ninja Multi-Config & MSVC / clang-cl & DLL
search path, import libraries, debug symbol PDBs \\
macOS arm64/x86-64 & Ninja / Xcode & AppleClang & bundle/framework
layout, signing and notarization \\
Android & Ninja & NDK Clang & ABI, SDK/NDK versions, app packaging \\
ARM Linux embedded & Ninja / Makefiles & cross GCC / Clang & sysroot,
libc and target architecture \\
\end{longtable}

This is a \emph{planning matrix}, not a statement that every target is
supported by the example or verified on the current machine. Define
minimum OS, architecture, compiler runtime and deployment constraints
for the real product.

\chapter{Configure model: native versus
cross-compiled}\label{configure-model-native-versus-cross-compiled}

A native build defaults to the host platform. A cross build must
represent a separate target platform before language/compilers are
enabled. The toolchain file provides target system, compiler, sysroot
and dependency-search controls. Avoid assigning compilers after the
\passthrough{\lstinline!project()!} call, since language detection has
already occurred.

\begin{lstlisting}
# toolchains/arm-linux-example.cmake (template; set real paths)
set(CMAKE_SYSTEM_NAME Linux)
set(CMAKE_SYSTEM_PROCESSOR aarch64)
set(CMAKE_C_COMPILER /opt/toolchains/bin/aarch64-linux-gnu-gcc)
set(CMAKE_CXX_COMPILER /opt/toolchains/bin/aarch64-linux-gnu-g++)
set(CMAKE_SYSROOT /opt/sysroots/arm64)
set(CMAKE_FIND_ROOT_PATH_MODE_PROGRAM NEVER)
set(CMAKE_FIND_ROOT_PATH_MODE_LIBRARY ONLY)
set(CMAKE_FIND_ROOT_PATH_MODE_INCLUDE ONLY)
set(CMAKE_FIND_ROOT_PATH_MODE_PACKAGE ONLY)
\end{lstlisting}

The snippet is a parameterized design pattern; it cannot run without a
matching toolchain/sysroot. Install prefixes and staging destinations
are separate concerns: a target sysroot is not necessarily the desired
package install root.

\chapter{Cross-compiler probe
strategy}\label{cross-compiler-probe-strategy}

Compiler checks may compile a program to confirm feature support.
\passthrough{\lstinline!try\_run!} requires running a binary, which
often fails in a cross environment without an emulator or preloaded
cache results. Use compile-only probes where possible; configure a
supported emulator only if execution semantics are trustworthy. Record
which feature checks were skipped and which were executed on target
hardware.

\section{Embedded-system acceptance
matrix}\label{embedded-system-acceptance-matrix}

\begin{longtable}[]{@{}lccc@{}}
\toprule\noalign{}
Verification & Host-only & Emulator & Actual target \\
\midrule\noalign{}
\endhead
\bottomrule\noalign{}
\endlastfoot
compile and link & Yes & Yes & Optional \\
ABI/layout assumptions & Partial & Partial & Yes \\
runtime loader behavior & No & Often & Yes \\
device drivers/peripherals & No & Simulated & Yes \\
real-time timing and energy & No & Inaccurate & Yes \\
\end{longtable}

A passing x86 build says little about ABI compatibility or device
peripherals on ARM. Preserve compiler ABI version, libc, sysroot and CPU
feature policy as release metadata.

\chapter{Windows runtime/deployment
engineering}\label{windows-runtimedeployment-engineering}

MSVC uses import libraries and DLLs in different roles. A library that
links successfully may still fail to launch because its DLL is not on
the runtime search path. CTest test-environment controls and runtime
dependency installation help, but must be paired with a verified release
layout. \passthrough{\lstinline!CMAKE\_MSVC\_RUNTIME\_LIBRARY!} requires
deliberate compatibility decisions when mixing static and dynamic
runtime usage; avoid assuming third-party libraries match the chosen
setting.

\chapter{Apple and Qt considerations}\label{apple-and-qt-considerations}

macOS frameworks and bundles require different installation and loader
strategies than typical Linux binaries. Validate
\passthrough{\lstinline!@rpath!}, embedded frameworks, architectures and
code signing at package testing time. Qt6 adds resources, translations,
platform plugins and deployment scripts; their presence is a runtime
contract, not proof of source compilation.

\chapter{Build performance model}\label{build-performance-model}

Let total elapsed engineering feedback time be:

\[
T_{feedback} = T_{configure} + T_{generate} + T_{compile} + T_{link} + T_{test} + T_{package}
\]

Measure each component separately. Optimization results are meaningful
only with a declared workload (clean build, header-change rebuild, one
\passthrough{\lstinline!.cpp!} edit, no-op build), identical hardware
and toolchain, and multiple independent samples.

\section{Main techniques and
trade-offs}\label{main-techniques-and-trade-offs}

\begin{longtable}[]{@{}
  >{\raggedright\arraybackslash}p{(\columnwidth - 6\tabcolsep) * \real{0.2500}}
  >{\raggedright\arraybackslash}p{(\columnwidth - 6\tabcolsep) * \real{0.2500}}
  >{\raggedright\arraybackslash}p{(\columnwidth - 6\tabcolsep) * \real{0.2500}}
  >{\raggedright\arraybackslash}p{(\columnwidth - 6\tabcolsep) * \real{0.2500}}@{}}
\toprule\noalign{}
\begin{minipage}[b]{\linewidth}\raggedright
Technique
\end{minipage} & \begin{minipage}[b]{\linewidth}\raggedright
Improves
\end{minipage} & \begin{minipage}[b]{\linewidth}\raggedright
Cost / risk
\end{minipage} & \begin{minipage}[b]{\linewidth}\raggedright
Acceptance experiment
\end{minipage} \\
\midrule\noalign{}
\endhead
\bottomrule\noalign{}
\endlastfoot
Ninja generator & scheduling overhead, incremental granularity & not all
IDE features & clean + incremental build timing \\
compiler cache (ccache/sccache) & repeated compilation & cache size and
key correctness & cold/warm cache hit ratio and wall time \\
precompiled headers & header parse time & compiler compatibility,
invalidation & common-header-change rebuild \\
unity builds & compile overhead across many files & symbol/name
collisions, diagnostics quality & clean build and error trace
comparison \\
parallelism & utilization of CPU & RAM pressure and IO contention &
controlled jobs-vs-time sweep \\
optimized target edges & reduce needless rebuilds & maintenance of
precise dependencies & single-file rebuild fan-out \\
alternative linkers & link throughput & portability and debugging
variation & measured link-time regression test \\
split DWARF where supported & debug object/link cost & toolchain and
cache compatibility & disk footprint and debug symbol usability \\
\end{longtable}

\textbf{Rule:} never present a measured speed-up without stating
baseline, sample size, machine, workload and variability.

\chapter{Experimental benchmark
protocol}\label{experimental-benchmark-protocol}

\begin{enumerate}
\def\labelenumi{\arabic{enumi}.}
\tightlist
\item
  Pin source commit, compiler, linker, CMake, generator, storage medium
  and runner specifications.
\item
  Use separate build trees for each candidate; do not rely on in-place
  cache mutations.
\item
  Warm the filesystem where appropriate, but report cold and warm cases
  separately.
\item
  Run at least five measurements per case; report median and spread, not
  the single fastest run.
\item
  Collect elapsed time, CPU utilization, peak RSS, disk consumption,
  cache hit rate and test failures.
\item
  Compare correctness (build + test + package) before declaring any
  optimization accepted.
\item
  Repeat on a representative CI runner and workstation to avoid
  optimizing only one environment.
\end{enumerate}

\section{Improvement calculations}\label{improvement-calculations}

For baseline \(T_b\) and candidate \(T_c\):

\[
\mathrm{Speedup}=\frac{T_b}{T_c},\qquad
\mathrm{TimeSaved}(\%)=100\frac{T_b-T_c}{T_b}.
\]

Example: if a representative incremental build drops from 80 s to 50 s,
speedup is \(1.6\times\) and time saved is \(37.5\%\). This is an
illustrative calculation, not a benchmark result. Under Amdahl's law,
optimizing only compilation cannot eliminate fixed test/packaging time.

\chapter{Header-dependency fan-out}\label{header-dependency-fan-out}

Public headers are included in consumer compilations; an unnecessary
transitive include can invalidate a large proportion of a project after
a small edit. Consider forward declarations where valid, PIMPL for
stable ABI boundaries, modular interfaces and narrower public dependency
requirements. Do not adopt C++20 modules solely as a speed optimization
without checking compiler, CMake and build-generator support and
measuring actual effects.

\chapter{Cache and reproducibility
controls}\label{cache-and-reproducibility-controls}

Compiler caches may be configured by
\passthrough{\lstinline!CMAKE\_CXX\_COMPILER\_LAUNCHER!} /
\passthrough{\lstinline!CMAKE\_C\_COMPILER\_LAUNCHER!} or environment
presets, depending on organization policy. Cache validity depends on
exact command line, inputs, compiler identity and relevant environmental
state. Treat unexpected cache hits as a correctness incident, not merely
a performance issue.

\begin{lstlisting}
find_program(CCACHE_PROGRAM ccache)
if(CCACHE_PROGRAM AND PROJECT_IS_TOP_LEVEL)
  set(CMAKE_CXX_COMPILER_LAUNCHER "${CCACHE_PROGRAM}")
endif()
\end{lstlisting}

A top-level project should preserve developer control: do not force a
cache, linker, sanitizer or CPU-specific flag on projects that include
the library as a dependency.

\chapter{Platform release
verification}\label{platform-release-verification}

\begin{itemize}
\tightlist
\item[$\square$]
  Configure and build from clean binary directories on every supported
  runner.
\item[$\square$]
  Exercise Debug and Release; distinguish single/multi-config
  generators.
\item[$\square$]
  Check CPU architecture, runtime library and executable loader
  dependencies.
\item[$\square$]
  Run CTest under realistic test environment and resource settings.
\item[$\square$]
  Stage install to a non-default prefix and run clean consumer test.
\item[$\square$]
  Inspect package layout, plugin dependencies and symbols (where
  relevant).
\item[$\square$]
  Run representative target-hardware verification for cross-compiled
  binaries.
\item[$\square$]
  Store toolchain version and benchmark metadata with release evidence.
\end{itemize}

\chapter{Performance decision record
template}\label{performance-decision-record-template}

\begin{longtable}[]{@{}
  >{\raggedright\arraybackslash}p{(\columnwidth - 2\tabcolsep) * \real{0.5000}}
  >{\raggedright\arraybackslash}p{(\columnwidth - 2\tabcolsep) * \real{0.5000}}@{}}
\toprule\noalign{}
\begin{minipage}[b]{\linewidth}\raggedright
Field
\end{minipage} & \begin{minipage}[b]{\linewidth}\raggedright
Entry
\end{minipage} \\
\midrule\noalign{}
\endhead
\bottomrule\noalign{}
\endlastfoot
Baseline / candidate & two immutable configurations \\
Workload & clean / no-op / source edit / header edit \\
Environment & CPU/RAM/disk/compiler/generator/version \\
Observations & samples, median, range, build/test result \\
Regression risks & memory use, debug experience, correctness, CI cost \\
Decision & accept / reject / trial with rollback trigger \\
\end{longtable}

\chapter{References}\label{references}

Textbook Chapters 24--26 and 43. CMake toolchain guide:
\url{https://cmake.org/cmake/help/latest/manual/cmake-toolchains.7.html}
; compile features:
\url{https://cmake.org/cmake/help/latest/manual/cmake-compile-features.7.html}
.

\end{document}
